\begin{figure}[htbp!]\centering
%		\resizebox{\linewidth}{!}{%
%			\begin{tikzpicture}[x=1.5cm,y=1.5cm,font=\scriptsize,
%				trave/.style={line width=1.7pt,line cap=round,line join=round,black},
%				biella/.style={line width=1.0pt,black},
%				carico/.style={-{Stealth[length=1.3mm]},line width=0.5pt,black!70},
%				reaz/.style={-{Stealth[length=2mm]},line width=0.9pt,gray!70},
%				bordo/.style={line width=0.5pt,black!70},
%				terra/.style={line width=0.6pt,black},
%				diagfill/.style={fill=gray!28,draw=black,line width=0.8pt},
%				nodoV/.style={circle,fill=white,draw=black,inner sep=0pt,minimum size=3.6pt,line width=0.6pt},
%				segno/.style={circle,draw=black,fill=white,inner sep=0.2pt,font=\tiny},
%				]
%				\def\sM{0.85}\def\sT{0.38}\def\sN{0.05}\def\sNc{0.18}
%				\newcommand{\nodi}{%
%					\coordinate (C) at (-0.5,0);\coordinate (D) at (1,0);\coordinate (E) at (2,0);
%					\coordinate (F) at (3,0);\coordinate (G) at (4.5,0);
%					\coordinate (A) at (0,-0.25);\coordinate (B) at (4,-0.25);}
%				\newcommand{\semicernD}[2]{\draw[black,fill=white,line width=0.9pt]($(#1)+(-0.025,0)$) arc[start angle=0,end angle={#2*180},radius=0.05]--cycle;}
%				\newcommand{\semicernF}[2]{\draw[black,fill=white,line width=0.9pt]($(#1)+(0.125,0)$) arc[start angle=0,end angle={#2*180},radius=0.05]--cycle;}
%				\newcommand{\fibre}{\draw[gray!55,dash pattern=on 2pt off 1.5pt,line width=0.45pt] (-0.5,-0.05)--(4.5,-0.05);}
%				\newcommand{\gnd}[1]{%
%					\draw[terra] ($(#1)+(-0.12,-0.18)$)--(#1)--($(#1)+(0.12,-0.18)$)--cycle;
%					\draw[terra] ($(#1)+(-0.18,-0.18)$)--($(#1)+(0.18,-0.18)$);
%					\foreach \t in {-0.14,-0.07,0,0.07,0.14}\draw[terra,line width=0.4pt] ($(#1)+(\t,-0.18)$)--++(-0.05,-0.07);}
%				\newcommand{\diagLin}[7]{%
%					\pgfmathsetmacro{\ax}{#5*#7*#3}\pgfmathsetmacro{\ay}{#5*#7*#4}%
%					\pgfmathsetmacro{\bx}{#6*#7*#3}\pgfmathsetmacro{\by}{#6*#7*#4}%
%					\filldraw[diagfill] (#1)--($(#1)+(\ax,\ay)$)--($(#2)+(\bx,\by)$)--(#2)--cycle;}
%				\newcommand{\diagPar}[8]{%
%					\pgfmathsetmacro{\ca}{2*(#5-2*#6+#7)}\pgfmathsetmacro{\cb}{-3*#5+4*#6-#7}\pgfmathsetmacro{\cc}{#5}%
%					\foreach \k in {0,1,...,20}{%
%						\pgfmathsetmacro{\tt}{\k/20}%
%						\pgfmathsetmacro{\mv}{(\ca*\tt*\tt+\cb*\tt+\cc)*#8}%
%						\pgfmathsetmacro{\ox}{\mv*#3}\pgfmathsetmacro{\oy}{\mv*#4}%
%						\coordinate (pp\k) at ($(#1)!\tt!(#2)+(\ox,\oy)$);}%
%					\filldraw[diagfill] (#1) \foreach \k in {0,1,...,20}{ -- (pp\k)} -- (#2) -- cycle;}
%				\newcommand{\assi}{%
%					\fibre
%					\draw[biella] (A)--(C);\draw[biella] (A)--(D);\draw[biella] (B)--(F);\draw[biella] (B)--(G);
%					\draw[trave] (C)--(G);
%					\gnd{A}\gnd{B}
%					\node[nodoV,xshift=2pt] at (C){};\node[nodoV] at (E){};\node[nodoV,xshift=-2pt] at (G){};
%					\node[nodoV] at (A){};\node[nodoV] at (B){};
%					\semicernD{D}{-1}\semicernF{F}{-1}}
%				
%				%==================== (a) STRUTTURA ====================
%				\begin{scope}[shift={(0,5.2)}]
%					\nodi
%					\draw[biella] (A)--(C);\draw[biella] (A)--(D);\draw[biella] (B)--(F);\draw[biella] (B)--(G);
%					\draw[trave] (C)--(G);
%					\draw[bordo] (1,0.225)--(3,0.225);
%					\foreach \x in {1,1.125,...,3.001}\draw[carico] (\x,0.225)--(\x,0);
%					\node[above] at (2,0.265){$p$};
%					\fibre \gnd{A}\gnd{B}
%					\draw[reaz] (A)--($(A)+(0.55,0)$);\node[gray!70,below,font=\tiny]at($(A)+(0.34,-0.04)$){$6p\ell$};
%					\draw[reaz] ($(A)+(0,-0.80)$)--($(A)+(0,-0.27)$);\node[gray!70,left,font=\tiny]at($(A)+(0.02,-0.62)$){$p\ell$};
%					\draw[reaz] (B)--($(B)+(-0.55,0)$);\node[gray!70,below,font=\tiny]at($(B)+(-0.34,-0.04)$){$6p\ell$};
%					\draw[reaz] ($(B)+(0,-0.80)$)--($(B)+(0,-0.27)$);\node[gray!70,right,font=\tiny]at($(B)+(-0.02,-0.62)$){$p\ell$};
%					\node[nodoV, xshift=2pt] at (C){};\node[nodoV] at (E){};\node[nodoV, xshift=-2pt] at (G){};
%					\semicernD{D}{-1}\semicernF{F}{-1}
%					\node[nodoV] at (A){};\node[nodoV] at (B){};
%					\node[left=-1pt] at (C){$C$};\node[below right=-0.2pt]at(D){$D$};\node[below=1pt] at (E){$E$};
%					\node[below left=0pt]at(F){$F$};\node[right=1pt] at (G){$G$};
%					\node[left=-1pt, yshift=-1pt] at (A){$A$};\node[right=1pt, yshift=-1pt] at (B){$B$};
%					\node[font=\tiny] at (2,-0.85){(a) schema, carico su $DF$ e reazioni};
%				\end{scope}
%				
%				%==================== (b) NORMALE N ====================
%				\begin{scope}[shift={(0,3.6)}]
%					\nodi
%					\diagLin{C}{D}{0}{1}{-0.66667}{-0.66667}{\sNc}
%					\diagLin{D}{E}{0}{1}{-6}{-6}{\sN}
%					\diagLin{E}{F}{0}{1}{-6}{-6}{\sN}
%					\diagLin{F}{G}{0}{1}{-0.66667}{-0.66667}{\sNc}
%					\assi
%					\node[segno] at ($(C)!0.5!(D)+(0,-0.55*0.66667*\sNc)$){$-$};
%					\node[segno] at ($(D)!0.5!(F)+(0,-0.5*6*\sN)$){$-$};
%					\node[segno] at ($(F)!0.5!(G)+(0,-0.55*0.66667*\sNc)$){$-$};
%					\node[below,font=\tiny] at ($(C)!0.5!(D)+(0,-0.66667*\sNc)$){$\tfrac23 p\ell$};
%					\node[below,font=\tiny] at ($(F)!0.5!(G)+(0,-0.66667*\sNc)$){$\tfrac23 p\ell$};
%					\node[below,font=\tiny] at ($(D)!0.5!(F)+(0,-6*\sN)$){$6 p\ell$};
%					\node[font=\tiny] at (2,-0.92){(b) sforzo normale $N$\ \ {\tiny($CD$,\,$FG$: ordinate amplificate)}};
%				\end{scope}
%				
%				%==================== (c) TAGLIO T ====================
%				\begin{scope}[shift={(0,1.9)}]
%					\nodi
%					\diagLin{C}{D}{0}{1}{-0.33333}{-0.33333}{\sT}
%					\diagLin{D}{E}{0}{1}{1}{0}{\sT}
%					\diagLin{E}{F}{0}{1}{0}{-1}{\sT}
%					\diagLin{F}{G}{0}{1}{0.33333}{0.33333}{\sT}
%					\assi
%					\node[segno] at ($(C)!0.5!(D)+(0,-0.2*\sT)$){$-$};
%					\node[segno] at ($(D)!0.45!(E)+(-0.15,0.30*\sT)$){$+$};
%					\node[segno] at ($(E)!0.55!(F)+(0.15,-0.30*\sT)$){$-$};
%					\node[segno] at ($(F)!0.5!(G)+(0,0.25*\sT)$){$+$};
%					\node[below,font=\tiny] at ($(C)!0.5!(D)+(0,-0.33333*\sT)$){$\tfrac13$};
%					\node[above,font=\tiny] at ($(D)+(0.06,1*\sT)$){$1$};
%					\node[below,font=\tiny] at ($(F)+(-0.06,-1*\sT)$){$1$};
%					\node[above,font=\tiny] at ($(F)!0.5!(G)+(0,0.33333*\sT)$){$\tfrac13$};
%					\node[font=\tiny] at (2,-0.72){(c) taglio $T$ \ ($\times\,p\ell$)};
%				\end{scope}
%				
%				%==================== (d) MOMENTO M ====================
%				\begin{scope}[shift={(0,0.4)}]
%					\nodi
%					\diagLin{C}{D}{0}{-1}{0}{-0.5}{\sM}
%					\diagPar{D}{E}{0}{-1}{-0.5}{-0.125}{0}{\sM}
%					\diagPar{E}{F}{0}{-1}{0}{-0.125}{-0.5}{\sM}
%					\diagLin{F}{G}{0}{-1}{-0.5}{0}{\sM}
%					\assi
%					\node[above,font=\tiny] at ($(D)+(0,0.5*\sM)$){$\tfrac12 p\ell^2$};
%					\node[above,font=\tiny] at ($(F)+(0,0.5*\sM)$){$\tfrac12 p\ell^2$};
%					\node[font=\tiny] at (2,-0.42){(d) momento flettente $M$};
%				\end{scope}
%		\end{tikzpicture}}
%		\caption{Caso (b), carico $p$ sul tratto centrale $DEF$ --- (a) schema e reazioni
%			($V_A=V_B=p\ell$, $H_A=H_B=6p\ell$); (b) sforzo normale $N$ (le bielle, scariche, portano $N$
%			uniforme: $N_{AC}=N_{BG}=+\tfrac{\sqrt5}{3}$, $N_{AD}=N_{BF}=-\tfrac{4\sqrt{17}}{3}$, in $p\ell$);
%			(c) taglio $T$; (d) momento flettente $M$, dal lato delle fibre tese. Scale indipendenti;
%			$M=0$ alle cerniere $C$, $E$, $G$ e ovunque non positivo (estradosso teso, lineare su $CD$, $FG$);
%			$N$ e $M$ simmetrici rispetto a $E$, $T$ emisimmetrico.}
%		\label{fig:stampella_NTM_b}
%	\end{figure}
